\section{Limitations}
\label{sec:limits}
We developed the compilation prompt on the development sets; the held-out test sets (\cref{sec:test}) were run once,
with gpt-6-luna only, and the baselines did not receive a comparable prompt search. gpt-6-luna is the
judge for every reader and also built every memory store, and it is the reader on which our only resolved gains appear,
so a preference of the judge for answers in its own style could inflate them; we have not yet measured agreement with
a second judge or with people. With three stronger readers our gains are small and mostly not resolved. The
development sets are small (200 questions in 10 conversations on BEAM-100K, 100 in 5 on BEAM-1M, 62 on
LongMemEval-S), BEAM conversations are synthetic, and LongMemEval-S is close to its ceiling for the systems we test. We
compare against one strong external system, Hindsight, and our re-implementation of LycheeMemory; other write-time
designs (Mem0, A-Mem, Zep) were not run under the same protocol, and we did not test hybrids that organize at both
write and read time. \method{} makes a compilation call for every question, which adds a second sequential call before the answer; it keeps every encoding error in the ledger; and it is weakest on summarization. The diagnoses in
\cref{sec:audit} rely on an LLM to decide whether information is present in a context; we checked them
against a second LLM but not against people.

\section{Conclusion}
\label{sec:conclusion}
We asked whether a long-term conversational memory should be organized when it is written or when it is read. On BEAM,
under a protocol that fixes the reader, the answer prompt, and the judge, the answer prompt alone (up to \prelim{8.5}
points) and the context budget (up to \prelim{21}) open gaps as large as those between systems, and the one
long-context baseline loses more points to evidence its reader fails to use than to evidence its retrieval fails to
find. Keeping the stored records unchanged and organizing them per question, as \method{} does, outperformed the
write-time system we compared against with gpt-6-luna as reader (\prelim{$+7.3$} and \prelim{$+4.2$} points, three runs);
with three stronger readers it was ahead by 1--3 points in every case, a gap resolved only once. These
results are on development sets and use gpt-6-luna as judge. Two observations may matter beyond our method. The
storage format matters more once something organizes it for the reader: notes compiled from atomic records beat
notes compiled from Hindsight's context by \prelim{6.6} points on BEAM-100K point questions, although the two formats
are close when read directly. And across our three readers, the benefit of read-time organization shrank as the reader used long contexts better, suggesting that how much a memory should organize depends on its reader; confirming this on held-out sets and with a second judge is the main open step.
